\section{Discussion and Limitations}

PRISM is designed as a deployable runtime security layer, not as a complete security boundary for LLM agents. Its primary contribution is to distribute enforcement across multiple runtime stages and to connect detection, policy enforcement, auditing, and operator workflows within a zero-fork gateway integration. That design provides practical value for agent gateways, but it also introduces a clear set of boundaries that should be stated explicitly.

\paragraph{Detection coverage is necessarily incomplete.} The heuristic layer is intentionally fast and operationally simple, which makes it suitable for inline use across multiple hooks, but it cannot cover all prompt-injection variants, obfuscation strategies, or context-dependent malicious instructions. The optional LLM-assisted scanner broadens coverage for suspicious tool-returned content, yet it remains a best-effort classifier rather than a formally trustworthy arbiter. Adaptive or novel adversarial inputs may evade both tiers. PRISM should therefore be understood as a layered runtime defense that raises the difficulty of successful attacks and narrows the exposure window, not as a mechanism that can guarantee perfect detection of all adversarial inputs.

\paragraph{Safety properties depend in part on explicit runtime policy, not only on classification.} This is a strength, because many failures of tool-using agents are better described as policy failures than as text-understanding failures. It is also a limitation, because policy quality must be maintained over time. Allowlists, denylists, protected-path rules, domain tiers, and DLP patterns all require review as deployment needs evolve. PRISM therefore reduces operational security burden in some areas by centralizing controls, but it does not eliminate the need for ongoing policy maintenance. Operators who deploy PRISM still bear responsibility for the correctness and currency of their configured policy surface.

\paragraph{File and path protections are intentionally scoped.} The current path-canonicalization logic performs string-level normalization and resolution, but it does not resolve symlink aliases through \texttt{realpath}. As a result, the system can enforce practical path checks for the common cases it models directly, but it should not be described as a complete filesystem sandbox. Likewise, the file-integrity monitor and audit chain are designed to make tampering visible and verifiable after the fact; they provide tamper-evident behavior, not tamper-proof storage. An attacker who can bypass string-level path resolution or who operates below the filesystem abstraction layer falls outside PRISM's protection boundary.

\paragraph{PRISM is framework-specific by design.} The zero-fork integration strategy depends on the hook surface and deployment assumptions of OpenClaw-based gateways. This specificity is part of why the implementation can be concrete and deployable, but it also limits direct portability. Applying the same architectural approach to other agent frameworks would likely preserve the overall defense model while requiring substantial re-mapping of lifecycle stages, tool abstractions, policy insertion points, and operational interfaces. The degree to which PRISM's mechanisms generalize beyond OpenClaw remains an open question that can only be answered through future porting efforts.

\paragraph{PRISM does not replace lower-layer system hardening.} It does not provide OS-level sandboxing, kernel isolation, hardware-rooted integrity, or general outbound network control outside the paths it directly mediates. A secure deployment should therefore treat PRISM as one layer in a broader defense stack that may also include host hardening, container or VM isolation, network egress controls, secret management infrastructure, and conventional service authentication. The absence of any of these complementary layers may leave attack surfaces that PRISM, by design, does not address.

\paragraph{The current system should not be overinterpreted as production-proven at large scale.} The implementation is real, runnable, and test-covered, and the evaluation methodology is designed to measure effectiveness, false positives, layer contribution, overhead, and operational recovery. Even so, large-scale multi-tenant validation, long-duration operational studies, and broader benchmark coverage remain future work. Claims about deployability are therefore strongest at the level of architecture and implementation maturity, not at the level of broad field validation.

\subsection{Future Directions}

These limitations also point toward several natural extensions. A productive next step would be to expand benchmark coverage with larger indirect-injection and tool-abuse corpora, then evaluate the system under more diverse workloads and deployment settings. A second direction is to combine PRISM's runtime hooks and policy layer with more formal safe-tool-use specifications, allowing deployable enforcement to be paired with stronger pre-deployment reasoning about tool constraints. Porting the architecture to other agent frameworks would help separate what is intrinsic to the PRISM design from what is specific to OpenClaw's lifecycle model. Finally, audit-driven policy refinement may provide a practical path toward improving rules and thresholds based on observed attack and false-positive patterns over time, closing the loop between runtime telemetry and policy evolution.

Overall, the appropriate interpretation of PRISM is not that it solves agent security in full, but that it demonstrates a practical systems approach for advancing agent defenses from isolated filters toward a lifecycle-aware, policy-enforced, and auditable runtime layer.
